\ifdefined\gkver\else\def\gkver{student}\fi
\documentclass[\gkver]{gaokaozhenti}
\begin{document}

\chapter{1978年全国}

\begin{enumerate}
\item
%% number: 1
%% paperName: 1978年全国高考第 1 题 2 分
%% typeId: 3
%% type: 填空题
%% chapter: 电磁感应
%% point: 法拉第电磁感应定律
%% method: 无
%% score: 2
%% degree: 850
%% duplicateId: 0
%% body:
当穿过一个线圈的 \tkanswer{磁通量} 发生变化时，线圈中产生感应电动势；感应电动势的大小，除与线圈的匝数成正比外，还与 \tkanswer{磁通量的变化率} 成正比。

%% answer: 磁通量；磁通量的变化率
\jdanswer{
磁通量；磁通量的变化率
}

%% memo:
\memoanswer{本题原卷及材料中未提供详解，待补充。}

\item
%% number: 2
%% paperName: 1978年全国高考第 2 题 2 分
%% typeId: 3
%% type: 填空题
%% chapter: 机械振动
%% point: 单摆
%% method: 无
%% score: 2
%% degree: 850
%% duplicateId: 0
%% body:
单摆在摆动过程中，其速度和加速度都是随时间变化的。从最大位移处向平衡位置运动的过程中，速度越来越 \tkanswer{大}，加速度越来越 \tkanswer{小}。

%% answer: 大；小
\jdanswer{
大；小
}

%% memo:
\memoanswer{本题原卷及材料中未提供详解，待补充。}

\item
%% number: 3
%% paperName: 1978年全国高考第 3 题 2 分
%% typeId: 3
%% type: 填空题
%% chapter: 原子和原子核
%% point: 天然放射性
%% method: 无
%% score: 2
%% degree: 850
%% duplicateId: 0
%% body:
在天然放射性元素的放射线中，已经查明，$\alpha$ 射线是 \tkanswer{氦核流或 $\ce{^{4}_{2}He}$ 流}，$\beta$ 射线是 \tkanswer{电子流}，$\gamma$ 射线是 \tkanswer{高频电磁波或光子流}。

%% answer: 氦核流或 ^{4}_{2}He 流；电子流；高频电磁波或光子流
\jdanswer{
氦核流或 $\ce{^{4}_{2}He}$ 流；电子流；高频电磁波或光子流
}

%% memo:
\memoanswer{本题原卷及材料中未提供详解，待补充。}

\item
%% number: 4
%% paperName: 1978年全国高考第 4 题 2 分
%% typeId: 3
%% type: 填空题
%% chapter: 机械波
%% point: 波长频率波速
%% method: 无
%% score: 2
%% degree: 950
%% duplicateId: 0
%% body:
在 $20\UdC$ 的空气中，声音的传播速度是 $340\Ums$。如果它的频率是 $100\UHz$，那么它的波长是 \tkanswer{$3.4\Um$}。

%% answer: 3.4 m
\jdanswer{
$3.4\Um$
}

%% memo:
\memoanswer{
由题意知 $f=100\UHz$，又 $T=\frac{1}{f}=0.01\Us$，由波速公式得 $v=\frac{\lambda}{T}$，故 $\lambda=vT=340\times0.01\Um=3.4\Um$。

评分标准：全题 10 分，每小题 2 分。
}

\item
%% number: 5
%% paperName: 1978年全国高考第 5 题 2 分
%% typeId: 3
%% type: 填空题
%% chapter: 电场
%% point: 库仑定律
%% method: 无
%% score: 2
%% degree: 950
%% duplicateId: 0
%% body:
两个点电荷之间距离为 $a$，相互作用力为 $f$；如果距离变为 $2a$，则相互作用力变为 \tkanswer{$\frac{f}{4}$}。

%% answer: \frac{f}{4}
\jdanswer{
$\frac{f}{4}$
}

%% memo:
\memoanswer{
由库仑定律得 $F=k\frac{Q_{1}Q_{2}}{a^{2}}=f$，故 $F'=k\frac{Q_{1}Q_{2}}{(2a)^{2}}=\frac{f}{4}$。

评分标准：全题 10 分，每小题 2 分。
}

\item
%% number: 6
%% paperName: 1978年全国高考第 6 题 10 分
%% typeId: 6
%% type: 计算题
%% chapter: 电路
%% point: 闭合电路欧姆定律
%% method: 无
%% score: 10
%% degree: 650
%% duplicateId: 0
%% body:
如图所示的电路中，三个电阻的阻值分别是 $R_{1}=2\UO$，$R_{2}=4\UO$，$R_{3}=4\UO$。电池的电动势 $E=4.2\UV$，内阻 $r=0.2\UO$。求：

\begin{enumerate}
	\item
	接通开关 K，断开开关 $K'$ 时，$R_{1}$ 和 $R_{2}$ 两端电压之比 $\frac{U_{1}}{U_{2}}$；
	\item
	两个开关都接通时，$R_{1}$ 和 $R_{2}$ 两端电压之比 $\frac{U_{1}'}{U_{2}'}$；
	\item
	两个开关都接通时，通过 $R_{1}$ 的电流强度 $I_{1}$。
\end{enumerate}
\onepicture[align=right, w=6cm]{figs/06.png}

%% answer: （1）\frac{1}{2}；（2）1；（3）1 A
\jdanswer{
\begin{enumerate}
	\item
	$\frac{1}{2}$

	\item
	1

	\item
	$1\UA$
\end{enumerate}
}

%% memo:
\memoanswer{
（1）K 接通、K$'$ 断开时，$R_{1}$ 与 $R_{2}$ 串联，通过两电阻的电流大小相等，由串联电路的分压原理得

$\frac{U_{1}}{R_{1}}=\frac{U_{2}}{R_{2}}$，解得 $\frac{U_{1}}{U_{2}}=\frac{R_{1}}{R_{2}}=\frac{2}{4}=\frac{1}{2}$。

（2）两个开关都接通时，$R_{2}$ 与 $R_{3}$ 并联，设并联电阻为 $R'$，则有

$R'=\frac{R_{2}R_{3}}{R_{2}+R_{3}}=\frac{4\times4}{4+4}=2\UO$，

$R'$ 与 $R_{1}$ 串联，由串联电路的分压原理得

$\frac{U_{1}'}{U_{2}'}=\frac{R_{1}}{R'}=\frac{2}{2}=1$。

（3）当两个开关都接通时，通过 $R_{1}$ 的电流，即干路电流，由闭合电路的欧姆定律得

$I=\frac{E}{R_{1}+R'+r}=\frac{4.2}{2+2+0.2}=1\UA$。

评分标准：全题 10 分，（1）3 分，（2）3 分，（3）4 分。

（1）中，直接写出 $\frac{U_{1}}{U_{2}}=\frac{R_{1}}{R_{2}}=\frac{2}{4}=\frac{1}{2}$ 的，同样给分。

（3）中，公式中漏了内阻 $r$ 的，不给分。
}

\item
%% number: 7
%% paperName: 1978年全国高考第 7 题 13 分
%% typeId: 6
%% type: 计算题
%% chapter: 光学
%% point: 透镜
%% method: 无
%% score: 13
%% degree: 650
%% duplicateId: 0
%% body:
用照相机对着一个物体照相，已知镜头（相当于一个凸透镜）的焦距为 $13.5\Ucm$，当底片与镜头的距离为 $15\Ucm$ 时，在底片上成 $5\Ucm$ 高的像。

\begin{enumerate}
	\item
	求物体的高；
	\item
	绘出光路图。
\end{enumerate}

%% answer: （1）45 cm；（2）光路图如图所示
\jdanswer{
\begin{enumerate}
	\item
	$45\Ucm$

	\item
	光路图如图所示
\end{enumerate}
}

%% memo:
\memoanswer{
（1）设物体高为 $H$，像高 $h=5\Ucm$，由透镜成像公式和放大率的定义得

$\frac{1}{u}+\frac{1}{v}=\frac{1}{f}$，即 $\frac{1}{u}+\frac{1}{15}=\frac{1}{13.5}$，

$\frac{h}{H}=\frac{v}{u}$，即 $\frac{h}{H}=\frac{15}{u}$，

联立解得 $H=45\Ucm$，$u=135\Ucm$，故物体的高度为 $45\Ucm$。

（2）光路图如图所示。

\onepicture[w=6cm]{figs/07_an.png}

评分标准：全题 13 分，（1）8 分，（2）5 分。

（1）中，能列出两个文字方程的给 2 分；经过文字运算，得到物体高度的文字表达式，并正确代入需要的数值的，再给 4 分；最后算出正确答案的，共给 8 分。列出两个文字方程后，立即正确代入需要的数值的，给 6 分；正确算出答案的，共给 8 分。

（2）中，用透镜前面或后面的焦点都可以，同样给分。图中未用箭头表示光线进行方向的，扣 1 分；图中尺寸完全不合比例的，扣 1 分。
}

\item
%% number: 8
%% paperName: 1978年全国高考第 8 题 13 分
%% typeId: 6
%% type: 计算题
%% chapter: 电路
%% point: 电表改装
%% method: 无
%% score: 13
%% degree: 650
%% duplicateId: 0
%% body:
一个电流$-$电压两用表的电路如图所示，电流计 G 的量程是 $0.001\UA$，内阻是 $100\UO$。两个电阻的阻值是 $R_{1}=9900\UO$，$R_{2}=1.01\UO$。

\begin{enumerate}
	\item
	双刀双掷开关接到哪边是电流表，接到哪边是电压表？
	\item
	电流表、电压表的量程各是多大？
\end{enumerate}
\onepicture[align=right, w=5cm]{figs/08.png}

%% answer: （1）见解析；（2）0.1 A，10 V
\jdanswer{
\begin{enumerate}
	\item
	双刀双掷开关接到 c、d 上，$R_{2}$ 与电流计 G 并联，组成电流表；若双刀双掷开关接到 a、b 上，电阻 $R_{1}$ 与电流计 G 串联，组成电压表。

	\item
	$0.1\UA$，$10\UV$
\end{enumerate}
}

%% memo:
\memoanswer{
（1）双刀双掷开关接到 c、d 上，$R_{2}$ 与电流计 G 并联，组成电流表；若双刀双掷开关接到 a、b 上，电阻 $R_{1}$ 与电流计 G 串联，组成电压表。

（2）分析电流表的量程 $I_{0}$。设干路电流为 $I$，通过电流计的电流为 $I_{g}$，由并联电路的特点得

$I_{g}r=(I-I_{g})R_{2}$，

解得 $I=\frac{R_{2}+r}{R_{2}}I_{g}=\frac{1.01+100}{1.01}\times0.001\UA\approx0.1\UA$。

分析电压表的量程 $U_{0}$。设电压表量程为 $U$，由串联电路的特点得

$U=I_{g}(r+R_{1})=0.001\times(100+9900)\UV=10\UV$。

评分标准：全题 13 分，（1）5 分；（2）两表的量程各 4 分，共 8 分。
}

\item
%% number: 9
%% paperName: 1978年全国高考第 9 题 14 分
%% typeId: 6
%% type: 计算题
%% chapter: 力物体的平衡
%% point: 浮力
%% method: 无
%% score: 14
%% degree: 650
%% duplicateId: 0
%% body:
一个 $14\Ug$ 重的密度计（如图所示），放在水中，水面在它的刻度 A 处；放在煤油中，油面在它的刻度 B 处。已知煤油的密度 $d=0.8\Ugcmc$，密度计刻度部分的玻璃管外半径 $r=0.75\Ucm$，求 AB 之间距离。
\onepicture[align=right, w=2cm]{figs/09.png}

%% answer: 2 cm
\jdanswer{
$2\Ucm$
}

%% memo:
\memoanswer{
设密度计刻度 A 以下的体积为 $V_{A}$，两种情况下浮力都等于密度计的重力，则有

水中：$F_{\text{浮}1}=\rho_{\text{水}}V_{\text{排}1}g$，$V_{\text{排}1}=V_{A}$；

煤油中：$F_{\text{浮}2}=\rho_{\text{油}}V_{\text{排}2}g$，$V_{\text{排}2}=V_{A}+\pi r^{2}\times\overline{AB}$；

又 $F_{\text{浮}1}=F_{\text{浮}2}=G$，联立解得 $AB\approx2\Ucm$。

评分标准：全题 14 分。能列出两个方程，表示没入液体中的体积乘液体比重等于比重计的重量的，给 7 分；进一步表示出两个体积之差同刻度 AB 的距离的关系的，再给 5 分；计算完毕无误的，给全分。答数写“$\approx2$ 厘米”“$=1.98$ 厘米”或“$=2$ 厘米”都作为正确答案。
}

\item
%% number: 10
%% paperName: 1978年全国高考第 10 题 20 分
%% typeId: 6
%% type: 计算题
%% chapter: 运动和力
%% point: 动量和能量
%% method: 无
%% score: 20
%% degree: 550
%% duplicateId: 0
%% body:
一质量 $M=2\Ukg$ 的木块，放在高 $h=0.8\Um$ 的光滑桌面上，被一个水平方向飞来的子弹打落在地面上（子弹留在木块中），落地点与桌边的水平距离 $s=1.6\Um$，子弹的质量 $m=10\Ug$。

\begin{enumerate}
	\item
	求子弹击中木块时的速度；
	\item
	子弹射入木块时产生的热量，若 90\% 被子弹吸收，子弹的温度能升高多少？（设子弹的比热为 $0.09$ 卡/克$\cdot$度，取 $g=10\Umsq$，空气阻力不计）
\end{enumerate}

%% answer: （1）804 m/s；（2）772 ℃
\jdanswer{
\begin{enumerate}
	\item
	$804\Ums$

	\item
	$772\UdC$
\end{enumerate}
}

%% memo:
\memoanswer{
（1）设子弹击中木块时的速度为 $v_{0}$，击中木块后的速度为 $v$，子弹击中木块的时间极短，由动量守恒得

$mv_{0}=(M+m)v$ ①

子弹击中木块后做平抛运动，由平抛规律得

$h=\frac{1}{2}gt^{2}$ ②

$s=vt$ ③

由②③联立解得 $v=s\sqrt{\frac{g}{2h}}=1.6\sqrt{\frac{10}{2\times0.8}}\Ums=4\Ums$ ④，

④代入①得

$v_{0}=\frac{M+m}{m}v=\frac{2+0.01}{0.01}\times4\Ums=804\Ums$。

（2）碰撞后，子弹和木块组成的系统损失的动能为 $\Delta E_{\mathrm{k}}$，由能量守恒得

$\Delta E_{\mathrm{k}}=\frac{1}{2}mv_{0}^{2}-\frac{1}{2}(m+M)v^{2}=\frac{1}{2}\times0.01\times804^{2}-\frac{1}{2}\times(0.01+2)\times4^{2}=3216\UJ$，

损失的动能 $\Delta E_{\mathrm{k}}$ 全部转换为热量 $Q$，则

$Q=3216\times0.24\Ucal\approx772\Ucal$，

又 $\eta Q=cm\Delta t$，

解得 $\Delta t=\frac{\eta Q}{cm}=\frac{90\%\times772}{0.09\times10}=772\UdC$，

故子弹升高的温度为 $772\UdC$。

评分标准：全题 20 分，（1）10 分，（2）10 分。

（1）中，列出动量守恒公式的，给 3 分；列出 $v=\frac{s}{t}$ 的，给 1 分；列出 $t=\sqrt{\frac{2h}{g}}$ 的，给 1 分；得出最后 $v$ 的文字式并正确代入数字的，再给 3 分；正确算出答案的，共给 10 分。如果分段进行数字计算，算出 $t$ 值的给 3 分，再算出 $v$ 值的再给 2 分，用动量守恒定律再算出 $v$ 的数值的再给 5 分，共 10 分。

（2）中，算出能量损耗数值的，给 4 分；换算成热量值的，再给 2 分；继续算出温度升高度数的，共给 10 分。如果把温度升高的度数说成子弹的温度，扣 3 分。
}

\item
%% number: 11
%% paperName: 1978年全国高考第 11 题 20 分
%% typeId: 6
%% type: 计算题
%% chapter: 电磁感应
%% point: 电磁感应中的变加速
%% method: 无
%% score: 20
%% degree: 550
%% duplicateId: 0
%% body:
如图所示，一个 U 形导体框架，宽度 $L=1\Um$，其所在平面与水平面交角 $\alpha=30^{\circ}$，其电阻可以忽略不计。设匀强磁场与 U 形框架的平面垂直，磁感应强度 $B=0.2\UT$。今有一条形导体 ab，其质量 $m=0.2\Ukg$，其有效电阻 $R=0.1\UO$，跨放在 U 形框上，并且能无摩擦地滑动。求：

\begin{enumerate}
	\item
	导体 ab 下滑的最大速度 $v_{m}$；
	\item
	在最大速度 $v_{m}$ 时，在 ab 上释放出来的电功率。
\end{enumerate}
\onepicture[align=right, w=6cm]{\includesvg{figs/11}}

%% answer: （1）2.5 m/s；（2）2.5 W
\jdanswer{
\begin{enumerate}
	\item
	$2.5\Ums$

	\item
	$2.5\UW$
\end{enumerate}
}

%% memo:
\memoanswer{
解法一：（1）导体 ab 在“下滑力” $mg\sin\alpha$ 的作用下沿 U 形框加速下滑切割磁感线，产生感应电动势为 $E=BLv$，感应电流的大小为 $I=\frac{E}{R}$，导体 ab 受到安培力的大小为 $F=BIL=\frac{B^{2}L^{2}v}{R}$。

随着切割速度增加，安培力也随之增大，导体 ab 的加速度逐渐减小；受力平衡时，导体 ab 的速度达到最大，由平衡条件得 $mg\sin\alpha=F$，解得 $v_{m}=\frac{mgR\sin\alpha}{B^{2}L^{2}}$，代入数据得 $v_{m}=2.5\Ums$。

（2）由电功率的表达式知 $P=IE=\frac{E^{2}}{R}=\frac{(BLv_{m})^{2}}{R}=2.5\UW$。

解法二：（1）导体 ab 匀速下滑时，重力势能转换为等量的电能，转换的即时功率值也应相等，显然在最大速度 $v_{m}$ 时有

$mg\sin\alpha\cdot v_{m}=\frac{E^{2}}{R}=\frac{(BLv_{m})^{2}}{R}$，

解得 $v_{m}=\frac{mgR\sin\alpha}{B^{2}L^{2}}=2.5\Ums$。

（2）导体 ab 上的电功率为 $P=\frac{E^{2}}{R}=mg\sin\alpha\cdot v_{m}=2.5\UW$。

评分标准：全题 20 分，（1）13 分，（2）7 分。

（1）中，列出感应电动势公式的给 2 分，列出安培力公式的再给 4 分，列出安培力同重力分量平衡的再给 5 分，算出 $v_{m}$ 的数值的共给 13 分。

（2）中，算功率无论从 $IE$、$I^{2}R$ 或 $\frac{E^{2}}{R}$ 进行计算，如果答案正确，都给 7 分。如果正确列出功率的详细文字公式而未代入数字算出答案的，给 5 分。
}

\end{enumerate}

\end{document}
